\section{总结与延伸}

本节把 EP102 压缩为七个结论。第一，CV 的 GPT 时刻不能简单复制 NLP，因为静态图像不天然携带人类语义和对齐。第二，多模态一体化的关键难点，是生成和理解无法互相校正。第三，next token prediction 是大模型起飞基础，也是推理跳步和精度不足的来源之一。第四，o1 的关键启发是用 RL 激发 CoT、反思和 critical decision 搜索。第五，视觉 Long CoT 是把慢思考迁移到视觉空间。第六，long context 不只是架构复杂度问题，更是压缩、检索、情景隔离和多模型协作问题。第七，下一个更深的 Agent 时刻，可能来自在线学习和 environmental scaling。

\begin{importantbox}{把 EP102 放进张小珺 AI/互联网队列}
EP102 是这一批里技术含量最高的多模态路线访谈之一。它连接 EP110 的 Agent 技术报告、EP113 的 agentic LLM、EP115 的 Agent 研究哲学，也为 EP106/EP109 的具身智能问题提供视觉推理和世界模型的上游解释。
\end{importantbox}

\subsection{术语总表}

前面各章已经分散解释了 VLM、CoT、动作空间、在线学习和世界模型。本节把这些术语集中放在一起，方便读者复盘整期论证链：从静态图像的数据属性，到 next token prediction 的目标函数，再到 o1 式反思、视觉 Long CoT 和在线学习。

\noindent\begin{tabularx}{\textwidth}{p{0.18\textwidth}p{0.36\textwidth}X}
\toprule
术语 & 简明定义 & 本期中的作用 \\
\midrule
VLM & Vision-Language Model，视觉语言模型 & 连接图像、视频和文本的主要模型形态。 \\
MIM & Masked Image Modeling，遮蔽图像建模 & CV 向 BERT 式预训练学习的代表路线。 \\
Next Token Prediction & 预测下一个 token 的训练目标 & 解释语言模型强大与跳步缺陷的共同来源。 \\
CoT & Chain of Thought，思维链 & 把单步决策展开成可搜索过程。 \\
Meta-CoT & 对 CoT 的再思考和反思 & 处理 critical decision 和路径选择。 \\
Action Space & 模型在任务中可选择的动作集合 & 决定 RL 是否能搜索到正确行为。 \\
Online Learning & 在线学习，系统在使用中持续更新 & 可能是下一代 Agent 的关键。 \\
World Model & 预测世界状态变化的内部模型 & 连接多模态、机器人和具身智能。 \\
\bottomrule
\end{tabularx}

\subsection{后续观察问题}

本节不是泛泛列问题，而是把访谈里的路线判断转成可追踪的观察点。接下来两年，如果多模态真的出现新的 GPT-4 时刻，应该能在这些位置看到信号：模型不只是 benchmark 上升，而是视觉推理、生成可控性、long context 和在线学习出现新的系统能力。

\begin{enumerate}
\item 视觉 Long CoT 能否在真实 benchmark 上明显提升多模态推理，而不只是生成更长解释？
\item 生成模型的可控性是否会通过视觉 CoT、局部编辑和反馈训练出现质变？
\item long context 会继续靠单模型扩窗，还是转向 planner/worker、多模型协作和工具化翻书？
\item 在线学习如何解决安全、数据污染、反馈稀疏和持续更新成本？
\item 机器人领域能否在抢跑硬件和本体之前，补上视觉推理、世界模型和动作反馈的基础能力？
\end{enumerate}

\subsection{拓展阅读}

\begin{itemize}
\item 对 Agent 训练和上下文工程感兴趣，可对照 EP110 Kimi K2 / ChatGPT Agent / Qwen3-Coder 技术报告笔记。
\item 对 Agentic LLM 和测试时扩展感兴趣，可对照 EP113 杨植麟 Kimi K2 访谈笔记。
\item 对具身智能的数据和世界模型问题感兴趣，可对照 EP106 王鹤、EP109 谢晨和 EP121 谭捷访谈。
\end{itemize}

\end{document}
